\documentclass[12pt]{article}
\usepackage{inputenc}
\usepackage[spanish]{babel}
\decimalpoint
\usepackage{geometry}
\usepackage{amsmath}
\usepackage{amssymb}
\usepackage{amsthm}
\usepackage{hyperref}
\usepackage{tikz}
\usetikzlibrary{babel}
\usepackage{titling}

\newgeometry{left=0.6in, right=0.6in, top=0.4in, bottom=0.5in}

\newtheorem{problem}{Problema}

\title{Procesos Estocásticos I $|$ Examen Final 2da Vuelta}
\author{Semestre 2026-2}
\date{}

\begin{document}

\maketitle
\vspace{-0.8cm}
\scriptsize

\textbf{Instrucciones generales:} Encierre sus respuestas finales en un recuadro. Justifique cada paso de sus cálculos y razonamientos.
% =============================================================

\begin{problem}
Considere una cadena de Markov $\{X_n : n \geq 0\}$ con espacio de estados $S = \{0, 1\}$. La distribución inicial
es $\pi = \!\left(\tfrac{1}{3},\, \tfrac{2}{3}\right)$ y la matriz de probabilidades de transición es:
\[
P = \begin{pmatrix}
\dfrac{1}{4} & \dfrac{3}{4} \\[8pt]
\dfrac{1}{2} & \dfrac{1}{2}
\end{pmatrix}.
\]
Calcule de forma explícita lo siguiente:
\begin{enumerate}

    \item[(a)] $\mathbb{P}(X_2 = 0,\; X_1 = 1,\; X_0 = 0)$.

    \item[(b)] $\mathbb{P}(X_{200} = 0 \mid X_{198} = 0)$.

    \item[(c)] $\mathbb{P}(X_2 = 1,\; X_0 = 1)$.
\end{enumerate}
\end{problem}

\begin{problem}
Un meteorólogo modela el clima de una ciudad con una cadena de Markov $\{X_n : n \geq 0\}$, donde
$n$ representa el día $n$-ésimo y el espacio de estados es $\mathcal{S} = \{\text{S}, \text{N}, \text{L}\}$
(Soleado, Nublado, Lluvioso). La matriz de transición diaria es:
\[
P = \begin{pmatrix}
1/2 & 1/2 & 0   \\[4pt]
1/4 & 1/2 & 1/4 \\[4pt]
1/2 & 0   & 1/2
\end{pmatrix},
\]
donde los estados están ordenados como S, N, L.
\begin{enumerate}
    \item[(a)] Argumente por qué esta cadena admite una única distribución estacionaria $\pi$ y encuéntrela.

    \item[(b)] Usando el Teorema Ergódico, estime cuántos días al año (de 365) serán en promedio
        soleados, nublados y lluviosos, respectivamente.

    \item[(c)] Sea $\tau_j = \inf\{n \geq 1 : X_n = j\}$ el tiempo de primer retorno al estado $j$
        y $m_{jj} = \mathbb{E}[\tau_j \mid X_0 = j]$ el tiempo medio de recurrencia. Use el Teorema Ergódico para calcular, cada cuántos días regresa en promedio un día lluvioso.
\end{enumerate}
\end{problem}



% =============================================================

\begin{problem}
Un algoritmo de recomendación de música clasifica las canciones en tres géneros:
\textbf{Pop~(P)}, \textbf{Rock~(R)} y \textbf{Jazz~(J)}. Después de reproducir una canción,
el algoritmo selecciona el género de la siguiente de acuerdo con la cadena de Markov
$\{X_n : n \geq 0\}$ con matriz de transición:
\[
P = \begin{pmatrix}
1/2 & 1/4 & 1/4 \\[4pt]
1/3 & 1/3 & 1/3 \\[4pt]
1/4 & 1/2 & 1/4
\end{pmatrix},
\]
donde los estados están ordenados como P, R, J.
Defina $\tau_A = \inf\{n \geq 1 : X_n = A\}$ como el primer tiempo de visita al estado $A$.
\begin{enumerate}
    \item[(a)] Un usuario está escuchando Jazz. Usando análisis de primer paso, calcule la probabilidad
        de que la siguiente canción de Pop llegue \textbf{antes} que la siguiente canción de Rock.
        Es decir, calcule $h_J = \mathbb{P}(\tau_P < \tau_R \mid X_0 = J)$.
        Plantee el sistema de ecuaciones de primer paso para $h_i$, $i \in \{P, R, J\}$, analice cuidadosamente las condiciones de frontera y resuélvalo.

    \item[(b)] Un usuario está escuchando Rock. Usando análisis de primer paso, calcule el número
        esperado de canciones que se reproducirán hasta escuchar Jazz por primera vez,
        es decir, $m_R = \mathbb{E}[\tau_J \mid X_0 = R]$.
        Plantee el sistema de ecuaciones de primer paso para $m_i$, $i \in \{P, R, J\}$, analice cuidadosamente las condiciones de frontera y resuélvalo.
\end{enumerate}
\end{problem}

% =============================================================
\begin{problem}
Las llamadas a un centro de atención médica llegan de acuerdo con un proceso de
Poisson $\{N_t\}_{t \geq 0}$ de tasa $\lambda = 4$ llamadas por hora. Sea $N_t$ el número de llamadas
recibidas en el intervalo $[0,t]$ (con $t$ en horas).

\begin{enumerate}

    \item[(a)] Calcule la probabilidad de que en los
        próximos 45 minutos llegue a lo más una llamada.

    \item[(b)] Cada llamada que llega al centro es clasificada de manera independiente como
        \textbf{urgente} con probabilidad $p = \tfrac{2}{3}$ o como \textbf{no urgente} con probabilidad
        $1-p = \tfrac{1}{3}$, sin importar las demás llamadas ni el instante de llegada.
        Sea $U_t$ el número de llamadas urgentes recibidas en $[0,t]$.
        Calcule $\mathbb{P}(U_t = k \mid N_t = n)$ para $k = 0, 1, \ldots, n$ e identifique la
        distribución resultante indicando sus parámetros.

\end{enumerate}
\end{problem}

\end{document}
